% GENERATED FILE. Edit canonical lessons, then run npm run generate:books.
% canonical-source-hash: fnv1a64:10fbc76a09572ac0
% canonical-lessons: AR-C456-hajm, AR-C456-sinf, AR-C456-ayyina, AR-C456-nasaba, AR-C456-daqa

\chapter{Finding the right size}
\label{ch:ar-finding-the-right-size}
\hlchaptermodality{456}

\begin{chapteropening}
\textbf{By the end of this chapter:} \emph{I can say volume, a variety, a sample, to fit and to be too tight.}

Complete the last lesson of chapter 456: I can say volume, a variety, a sample, to fit and to be too tight.
\end{chapteropening}

\section[\texorpdfstring{\ar{حجم} (ḥajm)}{حجم (ḥajm)}]{\ar{حجم} (ḥajm) — volume; size; bulk}

\label{lesson:AR-C456-hajm}



\begin{quote}
Before the new one: say the Arabic for to display, then the Arabic for a market stall.
\end{quote}

\subsection*{You'll want to know: \ar{حجم}}
\textbf{\ar{حجم}} — \emph{ḥajm} — \textquotedblleft{}volume; size; bulk\textquotedblright{}.

\begin{quote}
\ar{حجم} \ar{الحقيبة} \ar{صغير} \ar{جدا}.
\end{quote}

\begin{quote}
\emph{ḥajm al-ḥaqība ṣaghīr jiddan.}
\end{quote}

\begin{quote}
The size of the bag is very small.
\end{quote}

\emph{maqās} is a clothing size; \emph{ḥajm} is how big a thing is.

\begin{cousinweb}
\textbf{\ar{حجم}} (\emph{ḥajm}) is built on the native Arabic root \emph{ḥ-j-m}, \textquotedblleft{}the bulk of a thing\textquotedblright{} (Lane).
\end{cousinweb}

\begin{grammarlens}[title={What you've built}]
One more word for getting around.
\end{grammarlens}

\subsection*{Your turn}
\begin{itemize}
\raggedright
  \item \emph{Say it:} \emph{ḥajm}
  \item \emph{Say it:} \emph{ḥajm}, once more
  \item \emph{Say it:} \emph{ḥajm al-ḥaqība ṣaghīr jiddan.} — The size of the bag is very small.
  \item \emph{Recall:} say the Arabic for to display, then the Arabic for a market stall, then say \emph{ḥajm} again
\end{itemize}

\begin{tcolorbox}[breakable,colback=teal!4,colframe=teal!35!black,title={Before you move on}]
What does \ar{حجم} mean? (\textquotedblleft{}Volume; size; bulk\textquotedblright{}.) Say it once more.
\end{tcolorbox}
\section[\texorpdfstring{\ar{صنف} (ṣinf)}{صنف (ṣinf)}]{\ar{صنف} (ṣinf) — a variety; a kind}

\label{lesson:AR-C456-sinf}



\begin{quote}
Before the new one: say the Arabic for a market stall, then the Arabic for volume.
\end{quote}

\subsection*{You'll want to know: \ar{صنف}}
\textbf{\ar{صنف}} — \emph{ṣinf} — \textquotedblleft{}a variety; a kind\textquotedblright{}.

\begin{quote}
\ar{اشتريت} \ar{صنفا} \ar{جديدا} \ar{من} \ar{الشاي}.
\end{quote}

\begin{quote}
\emph{ishtaraytu ṣinfan jadīdan min ash-shāy.}
\end{quote}

\begin{quote}
I bought a new variety of tea.
\end{quote}

Plural: \ar{أصناف}, \emph{aṣnāf}.

\begin{cousinweb}
\textbf{\ar{صنف}} (\emph{ṣinf}) is built on the native Arabic root \emph{ṣ-n-f}, \textquotedblleft{}a sort, a kind\textquotedblright{} (Lane).
\end{cousinweb}

\begin{grammarlens}[title={What you've built}]
One more word for getting around.
\end{grammarlens}

\subsection*{Your turn}
\begin{itemize}
\raggedright
  \item \emph{Say it:} \emph{ṣinf}
  \item \emph{Say it:} \emph{ṣinf}, once more
  \item \emph{Say it:} \emph{ishtaraytu ṣinfan jadīdan min ash-shāy.} — I bought a new variety of tea.
  \item \emph{Recall:} say the Arabic for a market stall, then the Arabic for volume, then say \emph{ṣinf} again
\end{itemize}

\begin{tcolorbox}[breakable,colback=teal!4,colframe=teal!35!black,title={Before you move on}]
What does \ar{صنف} mean? (\textquotedblleft{}A variety; a kind\textquotedblright{}.) Say it once more.
\end{tcolorbox}
\section[\texorpdfstring{\ar{عينة} (ʿayyina)}{عينة (ʿayyina)}]{\ar{عينة} (ʿayyina) — a sample}

\label{lesson:AR-C456-ayyina}



\begin{quote}
Before the new one: say the Arabic for volume, then the Arabic for a variety.
\end{quote}

\subsection*{You'll want to know: \ar{عينة}}
\textbf{\ar{عينة}} — \emph{ʿayyina} — \textquotedblleft{}a sample\textquotedblright{}.

\begin{quote}
\ar{أعطاني} \ar{التاجر} \ar{عينة} \ar{صغيرة} \ar{من} \ar{العسل}.
\end{quote}

\begin{quote}
\emph{aʿṭānī at-tājir ʿayyina ṣaghīra min al-ʿasal.}
\end{quote}

\begin{quote}
The merchant gave me a small sample of the honey.
\end{quote}

\begin{cousinweb}
\textbf{\ar{عينة}} (\emph{ʿayyina}) is built on the native Arabic root \emph{ʿ-y-n}, \textquotedblleft{}the eye; a thing itself\textquotedblright{} (Lane). Modern Arabic uses it for a sample or a specimen.
\end{cousinweb}

\begin{grammarlens}[title={What you've built}]
One more word for getting around.
\end{grammarlens}

\subsection*{Your turn}
\begin{itemize}
\raggedright
  \item \emph{Say it:} \emph{ʿayyina}
  \item \emph{Say it:} \emph{ʿayyina}, once more
  \item \emph{Say it:} \emph{aʿṭānī at-tājir ʿayyina ṣaghīra min al-ʿasal.} — The merchant gave me a small sample of the honey.
  \item \emph{Recall:} say the Arabic for volume, then the Arabic for a variety, then say \emph{ʿayyina} again
\end{itemize}

\begin{tcolorbox}[breakable,colback=teal!4,colframe=teal!35!black,title={Before you move on}]
What does \ar{عينة} mean? (\textquotedblleft{}A sample\textquotedblright{}.) Say it once more.
\end{tcolorbox}
\section[\texorpdfstring{\ar{ناسب} (nāsaba)}{ناسب (nāsaba)}]{\ar{ناسب} (nāsaba) — to fit; to suit}

\label{lesson:AR-C456-nasaba}



\begin{quote}
Before the new one: say the Arabic for a variety, then the Arabic for a sample.
\end{quote}

\subsection*{You'll want to know: \ar{ناسب}}
\textbf{\ar{ناسب}} — \emph{nāsaba} — \textquotedblleft{}to fit; to suit\textquotedblright{}.

\begin{quote}
\ar{لم} \ar{يناسبني} \ar{القميص}.
\end{quote}

\begin{quote}
\emph{lam yunāsibnī al-qamīṣ.}
\end{quote}

\begin{quote}
The shirt did not suit me.
\end{quote}

Present: \ar{يناسب}, \emph{yunāsib}.

\begin{cousinweb}
\textbf{\ar{ناسب}} (\emph{nāsaba}) is built on the native Arabic root \emph{n-s-b}, \textquotedblleft{}kinship, relation\textquotedblright{} (Lane). It is Form III, \textquotedblleft{}to be in keeping with\textquotedblright{}; \emph{munāsib}, \textquotedblleft{}suitable\textquotedblright{}, is its participle.
\end{cousinweb}

\begin{grammarlens}[title={What you've built}]
One more word for getting around.
\end{grammarlens}

\subsection*{Your turn}
\begin{itemize}
\raggedright
  \item \emph{Say it:} \emph{nāsaba}
  \item \emph{Say it:} \emph{nāsaba}, once more
  \item \emph{Say it:} \emph{lam yunāsibnī al-qamīṣ.} — The shirt did not suit me.
  \item \emph{Recall:} say the Arabic for a variety, then the Arabic for a sample, then say \emph{nāsaba} again
\end{itemize}

\begin{tcolorbox}[breakable,colback=teal!4,colframe=teal!35!black,title={Before you move on}]
What does \ar{ناسب} mean? (\textquotedblleft{}To fit; to suit\textquotedblright{}.) Say it once more.
\end{tcolorbox}
\section[\texorpdfstring{\ar{ضاق} (ḍāqa)}{ضاق (ḍāqa)}]{\ar{ضاق} (ḍāqa) — to be too tight; to become narrow}

\label{lesson:AR-C456-daqa}



\begin{quote}
Before the new one: say the Arabic for a sample, then the Arabic for to fit.
\end{quote}

\subsection*{You'll want to know: \ar{ضاق}}
\textbf{\ar{ضاق}} — \emph{ḍāqa} — \textquotedblleft{}to be too tight; to become narrow\textquotedblright{}.

\begin{quote}
\ar{ضاق} \ar{الحذاء} \ar{على} \ar{قدمي}.
\end{quote}

\begin{quote}
\emph{ḍāqa al-ḥidhāʾ ʿalā qadamī.}
\end{quote}

\begin{quote}
The shoe was too tight on my foot.
\end{quote}

Present: \ar{يضيق}, \emph{yaḍīq}.

\begin{cousinweb}
\textbf{\ar{ضاق}} (\emph{ḍāqa}) is built on the native Arabic root \emph{ḍ-y-q}, \textquotedblleft{}to be narrow\textquotedblright{} (Lane). It is the plain verb beside \emph{ḍayyiq}, \textquotedblleft{}narrow\textquotedblright{}.
\end{cousinweb}

\begin{grammarlens}[title={What you've built}]
One more word for getting around.
\end{grammarlens}

\subsection*{Your turn}
\begin{itemize}
\raggedright
  \item \emph{Say it:} \emph{ḍāqa}
  \item \emph{Say it:} \emph{ḍāqa}, once more
  \item \emph{Say it:} \emph{ḍāqa al-ḥidhāʾ ʿalā qadamī.} — The shoe was too tight on my foot.
  \item \emph{Recall:} say the Arabic for a sample, then the Arabic for to fit, then say \emph{ḍāqa} again
\end{itemize}

\begin{tcolorbox}[breakable,colback=teal!4,colframe=teal!35!black,title={Before you move on}]
What does \ar{ضاق} mean? (\textquotedblleft{}To be too tight; to become narrow\textquotedblright{}.) Say it once more.
\end{tcolorbox}
